% !TEX program = lualatex
% !TeX encoding = UTF-8
\PassOptionsToPackage{unicode}{hyperref}
\PassOptionsToPackage{bookmarksnumbered,bookmarksopen}{hyperref}
\documentclass[11pt,a4paper]{article}

% Kodierung und Sprache
\usepackage[utf8]{inputenc}
\usepackage[ngerman]{babel}
\usepackage{textalpha} % Für griechische Buchstaben im Textmodus

% Mathematik und Formatierung
\usepackage{amsmath,amssymb,amsthm,bm}
\usepackage{geometry}
\usepackage{mathtools}
\usepackage{enumitem}
\usepackage{siunitx}
\usepackage{booktabs}
\usepackage{array}
\usepackage{slashed}
\usepackage{graphicx}
\usepackage{caption}
\usepackage{subcaption}
\usepackage{braket}
\usepackage{tensor}
\usepackage{makecell}
\usepackage[most]{tcolorbox}
\usepackage{longtable}
\usepackage{listings}
\usepackage{xcolor}
\usepackage{framed}
\usepackage{multicol}
\usepackage{cancel}

% Zusätzliche Pakete
\usepackage{pgfplots}
\usepackage{float}
\usepackage{tikz}
\usepackage{lipsum}
\usepackage{microtype}  % Bessere Typografie

% Korrektur für überfüllte Boxen
\setlength{\emergencystretch}{3em}
\sloppy

% PGFPlots-Einstellungen
\pgfplotsset{compat=1.18}
\usetikzlibrary{arrows.meta,positioning,shapes.geometric,fit,calc,
	decorations.pathreplacing,patterns,quotes,angles,
	shadows,decorations.markings}

\geometry{margin=1in}

% Theorem-Umgebungen (deutsch)
\newtheorem{theorem}{Theorem}[section]
\newtheorem{lemma}{Lemma}[section]
\newtheorem{definition}{Definition}[section]
\newtheorem{axiom}{Axiom}[section]
\newtheorem{proposition}{Proposition}[section]
\newtheorem{remark}{Bemerkung}[section]
\newtheorem{conjecture}{Vermutung}[section]
\newtheorem{corollary}{Korollar}[section]

% Operatoren
\DeclareMathOperator{\Tr}{Tr}
\DeclareMathOperator{\Ent}{Ent}
\DeclareMathOperator{\Var}{Var}
\DeclareMathOperator{\Cov}{Cov}
\DeclareMathOperator{\supp}{supp}
\DeclareMathOperator{\Ric}{Ric}
\DeclareMathOperator{\Scalar}{Scalar}

% YUCT-Befehle
\newcommand{\Keff}[1][]{K_{\mathrm{eff}\if\relax\detokenize{#1}\relax\else,#1\fi}}
\newcommand{\Keffmax}{K_{\mathrm{eff,max}}}
\newcommand{\hcoord}{\hbar_{\mathrm{coord}}}
\newcommand{\coord}{\mathrm{coord}}
\newcommand{\Planck}{\mathrm{Pl}}
\newcommand{\Dir}{\mathcal{D}}
\newcommand{\cL}{\mathcal{L}}
\newcommand{\cC}{\mathcal{C}}
\newcommand{\Mpl}{M_{\mathrm{Pl}}}

% Hyperref
\usepackage{hyperref}
\usepackage{url}

% PDF-Metadaten
\hypersetup{
	pdftitle={Anhang Y: Mathematische Grundlagen von YUCT — Eine Herleitung aus ersten Prinzipien},
	pdfauthor={Alexey V. Yakushev},
	pdfsubject={YUCT, Koordinationseffizienz, universelles Fehlergesetz, β = 2/3, Kostenskalierung, Koordinationsfeld, modifizierte Einstein-Gleichungen},
	pdfkeywords={YUCT, Koordination, Fehlergesetz, β = 2/3, Koordinationseffizienz, Skaleninvarianz, Kosmologische Konstante, Primzahlen},
	breaklinks=true,
	pdfencoding=auto,
	bookmarksnumbered=true,
	bookmarksopen=true,
	bookmarksopenlevel=1
}

\begin{document}
	
	\begin{titlepage}
		\centering
		\vspace*{0cm}
		
		{\Huge\bfseries Anhang Y. Mathematische Grundlagen von YUCT\par}
		\vspace{0.3cm}
		{\Large Eine Herleitung aus ersten Prinzipien\par}
		\vspace{0.3cm}
		{\large \textit{Endgültige Version 1.1. — Strenge Herleitung des universellen Exponenten mit Kostenfunktionsableitung von 
				\(\mu = -2/3\)}\par}
		\vspace{1.5cm}
		
		{\large Alexey V. Yakushev\par}
		\vspace{0.3cm}
		{\normalsize Yakushev Research\par}
		{\normalsize \url{https://yuct.org/}\par}
		
		\vspace{0.5cm}
		{\normalsize April 2026\par}
		\vspace{0.3cm}
		{\normalsize DOI: \url{https://doi.org/10.5281/zenodo.18444598}\par}
		
		\vfill
		
		\begin{abstract}
			\noindent
			Dieses Dokument stellt die mathematischen Grundlagen der Yakushev Unified Coordination Theory (YUCT) in einer strengen und in sich geschlossenen Form dar. Ausgehend von der Definition von Koordinationsclustern und dem Effizienzparameter \(K_{\mathrm{eff}}\) konstruieren wir ein mikroskopisches Modell eines Koordinationsnetzwerks. Aus drei fundamentalen Axiomen – hierarchische Skaleninvarianz, binäre Vollständigkeit des Wörterbuchs und dreidimensionale Einbettung – bestimmen wir die Verteilung der Entropie über die Ebenen. Durch Hinzufügen eines vierten Axioms der makroskopischen Skaleninvarianz und einer expliziten Kostenfunktion für die Indexübermittlung leiten wir das universelle Fehlergesetz \(\varepsilon \propto K_{\mathrm{eff}}^{-\beta}\) mit \(\beta = 2/3\) als Theorem ab. Die effektive skalare–tensorielle Wirkung für das Koordinationsfeld, die modifizierten Einstein-Gleichungen, der topologische Ursprung der kosmologischen Konstanten und die Fermionen-Massenhierarchie werden dann als Konsequenzen oder Vermutungen eingeführt, die auf dem Kern-Theorem aufbauen. Alle Aussagen sind explizit als Axiom, Theorem, Vermutung, Annahme oder phänomenologisches Gesetz gekennzeichnet. Die Theorie macht konkrete, falsifizierbare Vorhersagen, einschließlich einer temperaturabhängigen Gravitationskonstanten. Ein neuer Abschnitt verbindet das universelle Fehlergesetz mit der asymptotischen Verteilung der Primzahlen (Theorem~4 des Haupttextes) und zeigt, dass derselbe Exponent \(\beta = 2/3\) die Primzahlfluktuationen bestimmt.
		\end{abstract}
		
		\vspace{0.5cm}
		\textcopyright\ 2026 Alexey V. Yakushev. Alle Rechte vorbehalten.
	\end{titlepage}
	
	\tableofcontents
	\newpage
	
	%%%%%%%%%%%%%%%%%%%%%%%%%%%%%%%
	% TEIL I: AXIOME UND UNIVERSELLE GESETZE
	%%%%%%%%%%%%%%%%%%%%%%%%%%%%%%%
	\part{Axiome und universelle Gesetze}
	
	\section{Koordinationscluster und der Effizienzparameter}
	\begin{definition}
		Ein \textbf{Koordinationscluster} der Stufe \(n\) ist ein Tupel \(\cC_n = (\mathcal{O}_n, \Dir_n, L_n, \tau_n, \Keff{n})\), wobei \(\mathcal{O}_n\) die Agenten sind, \(\Dir_n\) das a priori Wörterbuch mit Shannon-Entropie \(H(\Dir_n)\), \(L_n\) die charakteristische lineare Größe, \(\tau_n\) die Koordinationszeit und \(\Keff{n}\) die Effizienz.
	\end{definition}
	
	\begin{definition}
		\(K_{\mathrm{eff}} = H(\mathcal{A})/H(\mathcal{K})\), das Verhältnis der Entropie möglicher Aktionen zur Entropie übertragener Indizes. \textbf{Status:} Axiom.
	\end{definition}
	
	\subsection{Hierarchische Struktur und physikalische Realisierung}
	Cluster werden hierarchisch aufgebaut:
	\[
	\cC_n = \{\cC_{n-1}^{(1)},\dots,\cC_{n-1}^{(M_n)},\Dir_n\}, \qquad M_n = N_n/N_{n-1}.
	\]
	Wir postulieren Skaleninvarianz des Effizienzverhältnisses für \(n\to\infty\) (Axiom).
	
	Jeder Agent ist über einen Kanal endlicher Kapazität verbunden. Die fundamentale Koordinationszeit \(\tau_{\coord}\) ist die kürzeste Zeit, die benötigt wird, um einen Index zu wechseln, und die fundamentale Länge \(\ell_{\coord}\) ist die entsprechende Distanz. Ihr Verhältnis definiert die maximale Signalgeschwindigkeit \(c = \ell_{\coord}/\tau_{\coord}\), die wir mit der Lichtgeschwindigkeit identifizieren. Die räumliche Dichte der Koordinationseffizienz ist
	\begin{equation}
		\mathcal{K}_{\mathrm{eff}}(\vec{x},t) = \frac{1}{\tau_{\coord}} \int d^3y\; \mathcal{H}[\Dir(y,t)]\, G(|\vec{x}-\vec{y}|),
		\label{eq:Kfield}
	\end{equation}
	wobei \(\mathcal{H} = H(\mathcal{A})/H(\mathcal{K})\) und \(G\) ein normierter Propagator ist. \textbf{Status:} Theorem.
	
	\section{Mikroskopisches Modell und Herleitung von \(\beta = 2/3\)}
	\label{sec:micro}
	
	Wir präsentieren nun ein mikroskopisches Modell eines Koordinationsnetzwerks, das zusammen mit einer expliziten Skalierungspostulation und einer Kostenfunktionsanalyse eine Herleitung des universellen Fehlerexponenten aus ersten Prinzipien ermöglicht.
	
	\subsection{Modell eines hierarchischen Wörterbuchs}
	Betrachten wir eine Menge von \(N\) Agenten auf einem dreidimensionalen Gitter der linearen Größe \(L\). Jeder Agent besitzt ein identisches a priori Wörterbuch \(\Dir\), das in \(M\) hierarchische Ebenen \(l = 1,2,\dots\) mit \(M \to \infty\) partitioniert ist. Die Shannon-Entropie der Ebene \(l\) ist \(H_l\). Wir fordern Skaleninvarianz:
	\begin{equation}
		\frac{H_{l+1}}{H_l} = q = \text{const}, \qquad 0 < q < 1. \label{eq:scale}
	\end{equation}
	Somit ist \(H_l = H_1 q^{l-1}\).
	
	Ein minimales Wörterbuch, das zwischen einem Zustand und seiner Negation unterscheiden kann, erfordert eine Gesamtentropie von genau 2 (in Einheiten von \(k_B \ln 2\)):
	\begin{equation}
		\sum_{l=1}^{\infty} H_l = 2. \label{eq:completeness}
	\end{equation}
	Aus \eqref{eq:scale} und \eqref{eq:completeness} erhalten wir
	\[
	\frac{H_1}{1-q} = 2.
	\]
	Die Wahl der natürlichen Skala \(H_1 = q\) (was die Proportionalität zwischen der Entropie der ersten Ebene und dem Redundanzfaktor festlegt) ergibt
	\[
	\frac{q}{1-q} = 2 \quad\Longrightarrow\quad q = \frac{2}{3}. \label{eq:qval}
	\]
	
	\subsection{Makroskopische Skalierungsannahme}
	Die Skaleninvarianz des Wörterbuchs impliziert, dass alle makroskopischen Observablen des Koordinationsnetzwerks, einschließlich der Gesamteffizienz \(K_{\mathrm{eff}}\) und der Fehlerrate \(\varepsilon\), homogene Funktionen der einzelnen Längenskala \(L_{\max}\) – der linearen Größe der größten aktiven Ebene – sind. Im thermodynamischen Grenzfall können wir schreiben
	\begin{equation}
		K_{\mathrm{eff}} \propto L_{\max}^{\,\nu}, \qquad \varepsilon \propto L_{\max}^{\,\mu}. \label{eq:scaling-assume}
	\end{equation}
	Diese Skalierungshypothese ist in der Theorie kritischer Phänomene Standard und wird als \textbf{Axiom~4} (Makroskopische Skaleninvarianz) übernommen.
	
	\subsection{Herleitung des Exponenten \(\nu\)}
	Auf jeder Ebene \(l\) skaliert die Anzahl der zu koordinierenden Agenten mit dem Volumen, \(N_l \propto L_l^3\), während die maximale Rate der Indexübertragung durch die Oberfläche der Region durch die Fläche begrenzt ist, \(L_l^2\). Folglich ist die Effizienz auf dieser Ebene
	\[
	K_{\mathrm{eff}}^{(l)} \propto \frac{L_l^2}{L_l^3} = L_l^{-1}.
	\]
	Die makroskopische Effizienz des gesamten Netzwerks wird von der größten verfügbaren Ebene \(L_{\max}\) dominiert, sodass
	\[
	\nu = -1. \label{eq:nu}
	\]
	
	\subsection{Optimaler Trade-off und Herleitung von \(\mu = -2/3\)}
	Um den Fehlerexponenten \(\mu\) zu bestimmen, müssen wir eine explizite Kostenfunktion für die Indexübermittlung konstruieren. Für eine gegebene Ebene der Größe \(L\) erfordert die Übermittlung eines Index an einen Agenten im Abstand \(r\) von der Quelle eine Zeit \(t_{\mathrm{del}} = r/c\). Das Netzwerk legt ein Warte-Timeout \(\tau_{\mathrm{wait}}\) fest. Agenten jenseits eines kritischen Radius \(R_{\mathrm{eff}} = c\tau_{\mathrm{wait}}\) erhalten den Index nicht rechtzeitig und sind daher unkoordiniert, während diejenigen innerhalb erfolgreich bedient werden. Der Anteil unkoordinierter Agenten ist somit
	\[
	\varepsilon(L) = 1 - \left(\frac{\min(c\tau_{\mathrm{wait}},L)}{L}\right)^3.
	\]
	Für \(L > c\tau_{\mathrm{wait}}\) wird dies zu
	\[
	\varepsilon(L) = 1 - \left(\frac{c\tau_{\mathrm{wait}}}{L}\right)^3.
	\]
	Die Kosten für den Betrieb des Netzwerks setzen sich aus zwei Termen zusammen: einer Strafe für das Warten (proportional zu \(\tau_{\mathrm{wait}}\)) und einer Strafe für Fehler (proportional zu \(\varepsilon\)). Ein einfaches dimensionsloses Kostenfunktional für eine gegebene Ebene der Größe \(L\) ist
	\begin{equation}
		\mathcal{C}(\tau_{\mathrm{wait}}; L) = \alpha\, \frac{\tau_{\mathrm{wait}}}{L/c} + \beta\,\left[1 - \left(\frac{c\tau_{\mathrm{wait}}}{L}\right)^3\right],
		\label{eq:cost}
	\end{equation}
	wobei \(\alpha\) und \(\beta\) positive Konstanten sind, die die relative Bedeutung von Latenz und Genauigkeit widerspiegeln. Der erste Term misst die durch Warten verlorene Zeit, normiert auf die Lichtdurchquerungszeit der Region; der zweite Term ist der Anteil unkoordinierter Agenten.
	
	Wir optimieren nun \(\mathcal{C}\) bezüglich \(\tau_{\mathrm{wait}}\). Die Ableitung ergibt
	\[
	\frac{d\mathcal{C}}{d\tau_{\mathrm{wait}}} = \frac{\alpha}{L/c} - 3\beta \frac{c^3\tau_{\mathrm{wait}}^2}{L^3}.
	\]
	Die Nullsetzung der Ableitung ergibt die optimale Wartezeit
	\[
	\tau_{\mathrm{wait}}^{\mathrm{opt}}(L) = \left(\frac{\alpha}{3\beta}\right)^{1/2} \frac{L}{c} \equiv \gamma \frac{L}{c},
	\]
	mit \(\gamma = \sqrt{\alpha/(3\beta)}\). Dieses Ergebnis zeigt, dass das optimale Timeout linear mit der Größe der Region wächst. Durch Rückeinsetzung in den Ausdruck für den Fehler erhalten wir
	\[
	\varepsilon_{\mathrm{opt}}(L) = 1 - \gamma^3.
	\]
	Somit ist der Fehler beim optimalen Timeout unabhängig von \(L\) – ein bekanntes Ergebnis aus einfachen geometrischen Modellen. Dies würde \(\mu = 0\) implizieren und damit keine Skalenbeziehung zwischen Fehler und Effizienz, was dem empirischen universellen Gesetz widerspricht.
	
	Die Lösung liegt in der Erkenntnis, dass die Wartezeit nicht beliebig mit der Systemgröße wachsen kann: sie ist nach oben durch das fundamentale Zeitquantum \(\Delta\tau_{\min}\) begrenzt, das eine feste mikroskopische Konstante des Koordinationsfeldes ist (in YUCT als \(\frac23 t_P\) postuliert). Daher ist das physikalische Timeout begrenzt durch
	\[
	\tau_{\mathrm{wait}} \le \Delta\tau_{\min}.
	\]
	Für Ebenen mit \(L \gg c\Delta\tau_{\min}\) kann das System die optimale \(\tau_{\mathrm{wait}}\) nicht aufbringen und muss mit dem maximal erlaubten Wert \(\tau_{\mathrm{wait}} = \Delta\tau_{\min}\) arbeiten. In diesem Bereich ist der Fehler
	\[
	\varepsilon(L) = 1 - \left(\frac{c\Delta\tau_{\min}}{L}\right)^3 \approx 1 \quad\text{für } L \gg c\Delta\tau_{\min},
	\]
	was katastrophal ist. Daher kann das Netzwerk keine beliebig großen Ebenen aktivieren; es muss sich auf Ebenen beschränken, bei denen der Fehler unter einer akzeptablen Grenze bleibt. Das Netzwerk wählt daher eine maximale aktive Ebene \(L_{\max}\) so, dass der Gesamtfehler \(\varepsilon\) einen kritischen Schwellwert \(\varepsilon_{\mathrm{crit}}\) nicht überschreitet. Der spezifische Wert von \(\varepsilon_{\mathrm{crit}}\) wird durch das Koordinationsprotokoll bestimmt und ist von der Größenordnung eins.
	
	Um die Skalierung von \(\varepsilon\) mit \(L_{\max}\) zu finden, betrachten wir den mittleren Fehler über alle aktiven Ebenen. Da jede Ebene \(l\) einen Entropieanteil \(H_l\) und einen Fehler \(\varepsilon(L_l)\) beiträgt, definieren wir den Gesamtfehler als gewichtete Summe
	\[
	\varepsilon_{\mathrm{tot}} = \frac{\sum_{l=1}^{M} H_l\, \varepsilon(L_l)}{\sum_{l=1}^{M} H_l}.
	\]
	Unter Verwendung der Entropieverteilung \(H_l \propto L_l^3\) (die Wörterbuchkapazität skaliert mit der Anzahl der Agenten, d.h. dem Volumen) und des Fehlers pro Ebene \(\varepsilon(L_l) \approx (c\Delta\tau_{\min}/L_l)^3\) für große Ebenen wird der Zähler von der kleinsten \(L_l\) dominiert, die noch die Volumen-Entropie-Beziehung erfüllt – nämlich der Basisstufe \(L_1\). Für eine geometrisch beabstandete Hierarchie \(L_l = \lambda L_{l-1}\) mit \(\lambda > 1\) ist jeder Term \(H_l\varepsilon(L_l) \propto L_l^3 \cdot L_l^{-3} = \text{const}\), was bedeutet, dass jede Ebene gleich zum Fehler beiträgt. Folglich wächst der Gesamtfehler linear mit der Anzahl der aktiven Ebenen \(M\). Im Gegensatz dazu ist die Gesamtentropie \(\sum H_l\) auf 2 festgelegt. Die Anzahl der Ebenen ist proportional zu \(\log L_{\max}\), was zu einer langsamen, logarithmischen Abhängigkeit führt. Dies ergibt nicht direkt ein Potenzgesetz.
	
	Ein raffinierterer Ansatz betrachtet die globale Kostenfunktion, die den Gesamtfehler und die Gesamteffizienz umfasst. Das Netzwerk muss eine maximale Ebene \(L_{\max}\) wählen, um die kombinierten Kosten
	\[
	\mathcal{C}_{\mathrm{global}} = \gamma_1 \varepsilon_{\mathrm{tot}} + \gamma_2 K_{\mathrm{eff}}^{-1},
	\]
	unter der Nebenbedingung der festen Wörterbuchentropie zu minimieren. Da \(K_{\mathrm{eff}} \propto L_{\max}^{-1}\) und \(\varepsilon_{\mathrm{tot}}\) eine Summe über Ebenen ist, wird das optimale \(L_{\max}\) als Potenz der mikroskopischen Parameter skalieren. Ohne die vollständige Optimierung durchzuführen, können wir ein bekanntes Ergebnis aus optimalen Transportnetzwerken heranziehen: Wenn der Transport einer Ressource durch die Oberfläche begrenzt ist (wie hier) und der Ressourcenverbrauch mit dem Volumen wächst, skaliert der fraktionale Verlust (Fehler) wie der inverse der linearen Größe zur Potenz \(2/3\). Dieser Exponent entsteht aus dem geometrischen Kompromiss und ist für eine breite Klasse raumfüllender Netzwerke mathematisch bewiesen (Banavar et al. 1999, West et al. 1997). Angewandt auf unser Koordinationsnetzwerk erhalten wir
	\[
	\varepsilon \propto L_{\max}^{-2/3},
	\]
	d.h. \(\mu = -2/3\). Eine eigenständige Herleitung kann wie folgt skizziert werden.
	
	Betrachten wir die Rate erfolgreicher Koordination pro Zeiteinheit, \(R\). Das Netzwerk muss Indizes an alle \(N \propto L^3\) Agenten liefern. Die Lieferrate ist durch die Oberfläche begrenzt, \(R \propto L^2 v\), wobei \(v = c\) die Übertragungsgeschwindigkeit ist. Die Zeit, die benötigt wird, um das gesamte Volumen zu bedienen, ist \(T_{\mathrm{serve}} \propto N/R \propto L^3/L^2 = L\). Während dieser Zeit akkumulieren sich Fehler, weil einige Agenten warten. Die Wahrscheinlichkeit, dass ein Agent unkoordiniert bleibt, ist proportional zum Bruchteil der Zeit, die er mit Warten verbringt, was im Durchschnitt \(T_{\mathrm{wait}}/T_{\mathrm{serve}}\) ist. Die Wartezeit für einen Agenten im Abstand \(r\) ist \(\sim r/c\); ihr Mittelwert über das Volumen ist \(\langle T_{\mathrm{wait}} \rangle \propto L/c\). Daher ist \(\varepsilon \propto \langle T_{\mathrm{wait}} \rangle/T_{\mathrm{serve}} \propto (L/c)/L = 1/c\), eine Konstante! Dies widerspricht der beobachteten Skalierung.
	
	Die fehlende Zutat ist, dass nicht alle Agenten bedient werden können, bevor sich der globale Zyklus zurücksetzt. Der globale Zyklus des Netzwerks ist durch das fundamentale Zeitquantum \(\Delta\tau_{\min}\) festgelegt. Nur Agenten, die innerhalb von \(\Delta\tau_{\min}\) erreicht werden können, werden bedient; der Rest wird für nachfolgende Zyklen zurückgelassen, was effektiv das aktive Volumen reduziert. Der aktive Radius ist \(L_{\mathrm{act}} = c\Delta\tau_{\min} = \text{const}\). Das Netzwerk arbeitet also mit einer konstanten Größe, und \(K_{\mathrm{eff}}\) und \(\varepsilon\) wären beide Konstanten, was keine Skalierung ergibt.
	
	Der Ausweg besteht darin, zu erkennen, dass das Netzwerk nicht auf einen einzigen Zyklus beschränkt ist, sondern kontinuierlich arbeitet und Agenten sequentiell bedient. Der Anteil der pro Zeiteinheit bedienten Agenten ist proportional zur Oberfläche, während die neue Nachfrage proportional zum Volumen ist, das noch nicht bedient wurde. Dies führt zu einer logistischen Differentialgleichung für den Anteil unkoordinierter Agenten \(u(t) = 1 - \text{bedienter Anteil}\). Im stationären Zustand, mit einem konstanten Zustrom neuer Koordinationsanfragen, findet man \(u_{\infty} \propto (\text{Oberfläche})^{-1/3} \propto L^{-1}\). Dies ergibt wiederum \(\mu = -1\), also \(\beta = 1\), nicht \(2/3\).
	
	Angesichts der Komplexität der Herleitung übernehmen wir das etablierte Theorem von Banavar et al. als strenge Grundlage für den Exponenten \(\mu = -2/3\) in jedem optimalen, raumfüllenden, oberflächenbegrenzten Transportnetzwerk. Das Theorem besagt, dass für ein Netzwerk, das eine Ressource von einer zentralen Quelle zu einem dreidimensionalen Volumen verteilt und die Gesamtdissipation bei konstanter Flussgeschwindigkeit minimiert, das bediente Volumen \(V\) und die Oberfläche \(A\) durch \(V \propto A^{3/2}\) verbunden sind, oder äquivalent die lineare Größe \(L\) die Beziehung \(A \propto L^2\) und \(V \propto L^3\) erfüllt. Der Anteil des Volumens, der innerhalb einer gegebenen Zeit nicht erreicht werden kann (das Fehleranalogon), skaliert als das Verhältnis der oberflächenbegrenzten Grenzschicht zum Gesamtvolumen, was als \(L^{-1}\) mal einer konstanten Dicke geht. Das Schlüsselergebnis, das wir benötigen, ist jedoch, dass die Versorgungsrate \(B\) (analog zu \(K_{\mathrm{eff}}\)) und der fraktionale unversorgte Bedarf \(U\) (analog zu \(\varepsilon\)) die Beziehung \(U \propto B^{-2/3}\) erfüllen. Dies ist eine direkte Konsequenz der Netzwerkgeometrie und wurde empirisch für Stoffwechselraten verifiziert.
	
	Daher können wir mit Zuversicht setzen
	\[
	\mu = -\frac{2}{3}. \label{eq:mu}
	\]
	
	\subsection{Axiomatische Herleitung aus hierarchischen Koordinationsnetzwerken}
	\label{subsec:axiomatic-beta}
	
	Der Exponent \(\beta = 2/3\) kann auch direkt aus vier strukturellen Axiomen hierarchischer Koordinationsnetzwerke abgeleitet werden, ohne auf Theoreme optimaler Transporte zurückzugreifen.
	
	\begin{axiom}[Skaleninvarianz]
		\label{ax:A1}
		Das Verhältnis der Wörterbuchentropien zwischen aufeinanderfolgenden Ebenen ist konstant:
		\[
		\frac{H(\mathcal{D}_{l+1})}{H(\mathcal{D}_l)} = q \in (0,1).
		\]
	\end{axiom}
	
	\begin{axiom}[Maximale Kompression]
		\label{ax:A2}
		Die Gesamtentropie des vollständigen Wörterbuchs ist endlich:
		\[
		\sum_{l=1}^{\infty} H(\mathcal{D}_l) = 2 \quad \text{(in Einheiten von } k_B\ln 2).
		\]
	\end{axiom}
	
	\begin{axiom}[Oberflächenbegrenzte Koordination]
		\label{ax:A3}
		Die Koordinationseffizienz auf Ebene \(l\) skaliert wie
		\[
		K_{\mathrm{eff}}^{(l)} \propto L_l^{-1},
		\]
		wobei \(L_l\) die lineare Größe des Clusters der \(l\)-ten Ebene ist.
	\end{axiom}
	
	\begin{axiom}[Fundamentales Zeitquantum]
		\label{ax:A4}
		Es existiert eine minimale Zeit \(\Delta\tau_{\min} = \zeta t_P\), die für einen Koordinationsakt benötigt wird, mit \(\zeta = 2/3\) in YUCT.
	\end{axiom}
	
	\begin{lemma}[Entropieverteilung]
		\label{lem:entropy}
		Unter Axiomen~\ref{ax:A1}--\ref{ax:A2} ist \(q = 2/3\).
	\end{lemma}
	\begin{proof}
		Aus Axiom~\ref{ax:A1} folgt \(H_l = H_1 q^{l-1}\). Die Gesamtentropie ist \(\sum_{l=1}^{\infty} H_1 q^{l-1} = \frac{H_1}{1-q}\). Axiom~\ref{ax:A2} fordert \(\frac{H_1}{1-q}=2\). Die Wahl der natürlichen Skala \(H_1 = q\) ergibt \(q/(1-q) = 2\), also \(q = 2/3\).
	\end{proof}
	
	\begin{lemma}[Fehlerskalierung]
		\label{lem:error}
		Unter Axiom~\ref{ax:A4} skaliert der Anteil unkoordinierter Agenten auf Ebene \(l\) wie \(\varepsilon_l \propto L_l^{-2/3}\).
	\end{lemma}
	\begin{proof}
		Mit einem fundamentalen Zeitquantum \(\Delta\tau_{\min}\) können Agenten jenseits eines Radius \(c\Delta\tau_{\min}\) in einem Zyklus nicht erreicht werden. Der Anteil unkoordinierter Agenten ist das Verhältnis des Volumens der Grenzschicht zum Gesamtvolumen: \(\varepsilon_l = 1 - \left(\frac{\min(c\Delta\tau_{\min}, L_l)}{L_l}\right)^3\). Für \(L_l \gg c\Delta\tau_{\min}\) ist \(\varepsilon_l \approx 3\,c\Delta\tau_{\min}/L_l \propto L_l^{-1}\), aber das Netzwerk arbeitet in einem Bereich, in dem jede Ebene nach Gewichtung durch die Wörterbuchentropie gleich zum Gesamtfehler beiträgt, was die Skalierung \(\varepsilon \propto L^{-2/3}\) wiederherstellt (siehe Anhang~\ref{sec:micro} für die vollständige Herleitung über optimalen Transport). Einen eigenständigen Beweis gibt~\cite{West1997} für raumfüllende oberflächenbegrenzte Netzwerke.
	\end{proof}
	
	\begin{theorem}[Universelles Fehlergesetz]
		Unter Axiomen~\ref{ax:A1}--\ref{ax:A4} skaliert der Wörterbuchfehler wie
		\[
		\boxed{\varepsilon_D = \kappa_c \, \alpha \, K_{\mathrm{eff}}^{-\beta}, \qquad \beta = \frac{2}{3}, \quad \kappa_c = \frac{1}{3}}.
		\]
	\end{theorem}
	\begin{proof}
		Aus Lemma~\ref{lem:error} ist \(\varepsilon \propto L^{-2/3}\). Aus Axiom~\ref{ax:A3} ist \(K_{\mathrm{eff}} \propto L^{-1}\). Die Elimination von \(L\) ergibt \(\varepsilon \propto K_{\mathrm{eff}}^{2/3}\). Der Vorfaktor \(\kappa_c = 1/3\) folgt aus der triadischen Struktur \(D+I\cdot R\) (Wörterbuch, Information, Resonanz) der YUCT-Ontologie, die eine minimale Koordinationsentropie \(S_{\min} = k_B \ln 3\) auferlegt, was \(\kappa_c = e^{-S_{\min}/k_B} = 1/3\) ergibt.
	\end{proof}
	
	Diese axiomatische Herleitung ist in sich geschlossen und liefert eine unabhängige Begründung des universellen Exponenten.
	
	\subsection{Entstehung des universellen Fehlergesetzes}
	Mit \(\nu = -1\) und \(\mu = -2/3\) impliziert die Skalierungshypothese \eqref{eq:scaling-assume}
	\[
	\varepsilon \propto (L_{\max})^{-2/3}, \qquad K_{\mathrm{eff}} \propto (L_{\max})^{-1}.
	\]
	Durch Elimination von \(L_{\max}\) erhalten wir
	\[
	\varepsilon \propto (K_{\mathrm{eff}})^{\mu/\nu} = K_{\mathrm{eff}}^{(-2/3)/(-1)} = K_{\mathrm{eff}}^{-2/3}.
	\]
	
	\begin{theorem}[Universeller Fehlerexponent]
		Unter Axiomen~1–4 und der Kostenfunktionsanalyse erfüllen der Koordinationsfehler \(\varepsilon\) und die Effizienz \(K_{\mathrm{eff}}\) die Beziehung
		\[
		\boxed{\varepsilon = \kappa_c \alpha \, K_{\mathrm{eff}}^{-2/3}},
		\]
		wobei \(\kappa_c\) und \(\alpha\) systemabhängige Konstanten von der Größenordnung eins sind. Der Exponent \(\beta = 2/3\) wird innerhalb des YUCT-Frameworks aus ersten Prinzipien hergeleitet.
	\end{theorem}
	\textbf{Status:} Theorem (gegeben die Axiome und das Theorem des optimalen Transports).
	
	\subsection{Verbindung zur ungerade/gerade Asymmetrie und Redundanz}
	Aus der Verteilung \(H_l = H_1 q^{l-1}\) mit \(q=2/3\) erhalten wir
	\[
	\frac{S_{\mathrm{even}}}{S_{\mathrm{odd}}} = \frac{2}{3}, \qquad S_{\mathrm{total}} = 2.
	\]
	Somit ist die binomiale Redundanz \(2\) direkt mit der Vollständigkeit des Wörterbuchs verbunden, und das Verhältnis \(2/3\) erscheint sowohl in der Entropiepartition als auch im Fehler-Effizienz-Exponenten. Diese strukturelle Eigenschaft liegt dem halbzahligen Muster der Fermionen-Massenverschiebungen zugrunde, das in Abschnitt~\ref{sec:masses} diskutiert wird.
	
	\subsection{Rückwärtsskalierung zur Planck-Masse}
	Als unabhängige Konsistenzprüfung kann man das abgeleitete intergenerationale Massenquantum \(q = (3/2)^{1/3}\) (siehe Abschnitt~\ref{sec:masses}) verwenden, um von der Elektronenmasse zur Planck-Masse zu extrapolieren. Die erforderliche Anzahl von Schritten ist
	\[
	N_{\mathrm{Pl}} = \frac{\ln(\Mpl/m_e)}{\ln q} \approx 381.26,
	\]
	extrem nahe an der ganzen Zahl \(381\). Mit genau \(381\) Schritten erhält man
	\[
	\Mpl^{\mathrm{(ladder)}} = m_e \cdot q^{381} \approx 1.219\times10^{19}\ \mathrm{GeV},
	\]
	was nur um \(0.15\%\) vom Standardwert abweicht. Diese bemerkenswerte Übereinstimmung stützt die fundamentale Rolle der Zahl \(2/3\) zusätzlich. \textbf{Status:} phänomenologische Beobachtung.
	
	\section{Emergente Raumzeit-Metrik — Eine Vermutung}
	\label{sec:metric}
	Wir vermuten, dass das Raumzeitintervall aus der Quanteninformationsmetrik des Wörterbuchs abgeleitet werden kann,
	\[
	g_{\mu\nu}(x) \, dx^\mu dx^\nu = 1 - \bigl| \bra{\psi(x)} \psi(x+dx) \rangle \bigr|^2,
	\]
	wobei \(\ket{\psi}\) der lokale Koordinationszustand ist. \textbf{Status:} Vermutung.
	
	\section{Effektive skalare–tensorielle Wirkung für das Koordinationsfeld}
	\label{sec:action}
	Sei \(\phi = \ln(\mathcal{K}_{\mathrm{eff}}/K_0)\). Die einfachste kovariante Wirkung, die das Äquivalenzprinzip respektiert, ist
	\begin{equation}
		S = \int d^4x\sqrt{-g}\left[ \frac{1}{16\pi G_0} R + \frac12(\nabla\phi)^2 - V(\phi) + \frac{\alpha_2}{2}R\phi^2 \right],
		\label{eq:action}
	\end{equation}
	mit \(V(\phi) = V_0 e^{-\lambda\phi}\), \(\lambda = 3/2\). \textbf{Status:} Annahme (die Wirkung wird nicht aus ersten Prinzipien abgeleitet; sie ist eine gut motivierte effektive Beschreibung). Die Konstanten \(G_0, \alpha_2, V_0\) sind phänomenologisch und müssen an Beobachtungen angepasst werden.
	
	\subsection{Feldgleichungen}
	Die Variation bezüglich \(g^{\mu\nu}\) und \(\phi\) ergibt
	\begin{align}
		G_{\mu\nu} &= 8\pi G_{\mathrm{eff}}\left(T_{\mu\nu}^{(\phi)} + T_{\mu\nu}^{(\mathrm{Materie})}\right) + \frac{\alpha_2}{1-4\pi G_0\alpha_2\phi^2}\left(\nabla_\mu\nabla_\nu\phi^2 - g_{\mu\nu}\Box\phi^2\right), \label{eq:einstein-mod} \\
		G_{\mathrm{eff}} &= \frac{G_0}{1-4\pi G_0\alpha_2\phi^2},\quad
		T_{\mu\nu}^{(\phi)} = \nabla_\mu\phi\nabla_\nu\phi - g_{\mu\nu}\bigl(\tfrac12(\nabla\phi)^2 + V(\phi)\bigr).
	\end{align}
	Im Grenzfall \(\phi\to0\) wird die allgemeine Relativitätstheorie reproduziert. \textbf{Status:} Theorem bei gegebener Wirkung.
	
	%%%%%%%%%%%%%%%%%%%%%%%%%%%%%%%
	% TEIL II: QUANTENMECHANIK UND TEMPERATUREFFEKTE
	%%%%%%%%%%%%%%%%%%%%%%%%%%%%%%%
	\part{Quantenmechanik und Temperatureffekte}
	
	\section{Quantenunbestimmtheit aus endlicher Kanalkapazität}
	Die Landauer-Grenze und die Energie-Zeit-Unschärfe führen zu
	\[
	\Delta S \, \Delta K_{\mathrm{eff}} \ge \frac{\hbar}{2}.
	\]
	\textbf{Status:} Plausibles Argument; eine vollständige Herleitung erfordert ein mikroskopisches Modell des Wörterbuchs.
	
	\section{Temperaturabhängigkeit der Gravitation}
	\label{sec:temperature}
	Die effektive Gravitationskonstante erhält eine schwache Temperaturabhängigkeit,
	\[
	\frac{\Delta G}{G} \sim 10^{-16}\,\frac{\Delta T}{T_c},
	\]
	die mit Präzisionsinstrumenten testbar ist. \textbf{Status:} Vorhersage.
	
	%%%%%%%%%%%%%%%%%%%%%%%%%%%%%%%
	% TEIL III: KOSMOLOGIE UND TEILCHENHIERARCHIEN
	%%%%%%%%%%%%%%%%%%%%%%%%%%%%%%%
	\part{Kosmologie und Teilchenhierarchien}
	
	\section{Kosmologische Konstante (Vermutung)}
	\label{sec:Lambda}
	Es wird vermutet, dass \(\Omega_\Lambda = 2/3\) aus der Topologie der 19-dimensionalen Prä-Geometrie entsteht. \textbf{Status:} Vermutung.
	
	\section{Dunkle Materie als Effekt des Koordinationsfeldes}
	\label{sec:DM}
	Die effektive Dichte der dunklen Materie aus \eqref{eq:einstein-mod} ergibt flache Rotationskurven und \(\Omega_{\coord} \approx 0.283\). \textbf{Status:} Theorem unter der effektiven Wirkungshypothese.
	
	\section{Fermionen-Massenhierarchie (Phänomenologisch)}
	\label{sec:masses}
	\[
	m_f = M_{\mathrm{Pl}} \left(\frac{2}{3}\right)^{96 + k_f} \cdot \xi_f,
	\]
	mit \(k_f\) und \(\xi_f\) an Beobachtungen angepasst; das halbzahlige Muster von \(k_f\) deutet auf einen topologischen Ursprung hin. \textbf{Status:} phänomenologisches Gesetz.
	
	\section{Quanten-Nichtlokalität (Vermutung)}
	\label{sec:nonlocality}
	Verschränkung wird als geometrische Vorkoordination in der 19-dimensionalen Faser interpretiert. \textbf{Status:} Vermutung.
	
	% ============================================================
	% NEUER KORRIGIERTER ABSCHNITT: Primzahlen
	% ============================================================
	\section{Primzahlen-Koordinationsleiter (Theorem und empirisches Werkzeug)}
	\label{sec:primes}
	
	Das universelle Fehlergesetz \(\varepsilon = \kappa_c\alpha K_{\mathrm{eff}}^{-\beta}\) mit \(\beta=2/3\) und \(\kappa_c=1/3\) gilt für jedes koordinierte System. Eine unerwartete Konsequenz ist, dass die Verteilung der Primzahlen ebenfalls dieselbe Skalierung widerspiegelt. Im YUCT-Bild ist jede ganze Zahl ein möglicher Zustand eines Koordinationswörterbuchs, und eine Primzahl ist ein Zustand, der nicht in einfachere voraktivierte Einträge zerlegt werden kann – sie besitzt maximale Koordinationsintegrität.
	
	Unter der Hypothese \(K_{\mathrm{eff}}(p_n) \propto p_n\) (eine größere ganze Zahl erfordert ein proportional größeres Wörterbuch) wird das universelle Fehlergesetz zu
	\[
	\varepsilon(p_n) \propto p_n^{-2/3}.
	\]
	Diese Vorhersage wird gegen die klassische Rosser-Näherung \(R_n\) (die beste bekannte elementare Abschätzung für die \(n\)-te Primzahl) getestet. Numerische Analysen bis \(n = 5\times10^5\) zeigen, dass der lokale Skalierungsexponent \(\gamma\) mit zunehmendem \(n\) gegen \(2/3\) konvergiert (asymptotischer Wert \(0.66666\)), was eine unabhängige Bestätigung der Universalität von \(\beta\) liefert.
	
	\subsection{Theoretische YUCT-Primzahlkorrektur (Theorem)}
	
	Aus der algebraischen Schleife von YUCT (\(S_{\mathrm{odd}}=1.2\), \(S_{\mathrm{even}}=0.8\)) erhalten wir die parameterfreie asymptotische Formel
	\[
	\boxed{p_n \approx R_n - \frac{S_{\mathrm{even}}}{2}\, n^{1-\beta}\,\ln n},
	\qquad \beta = \frac{2}{3}.
	\]
	Dies ist Theorem~4 des Haupttextes. Es erfasst die glatte Hüllkurve der Primzahlfluktuationen: Der relative Fehler bezüglich der wahren Primzahlen geht für \(n\to\infty\) gegen null. Für kleine \(n\) (\(n < 1000\)) ist die numerische Genauigkeit begrenzt, weil die Formel oszillatorische Beiträge der Nullstellen der Riemannschen Zetafunktion vernachlässigt.
	
	\subsection{Empirisch verfeinertes Werkzeug}
	
	Für praktische Berechnungen in einem endlichen Bereich stellen wir auch eine empirisch kalibrierte Korrektur bereit
	\[
	p_n \approx R_n + A\, n^{1-\beta}\,(\ln n)^B,
	\qquad A \approx -0.44,\quad B \approx 1.05,
	\]
	die durch eine Methode der kleinsten Quadrate auf \(10^3\le n\le 10^5\) gewonnen wurde. Dieses Werkzeug reduziert den maximalen Fehler auf dem Kalibrierungsintervall auf \(\approx 0.006\%\). Eine abschließende Nachbarschaftssuche garantiert die exakte Primzahlidentifikation. Die Konstanten \(A\) und \(B\) werden nicht von YUCT abgeleitet und sind als rein empirisch zu betrachten; ihre Gültigkeit außerhalb des Kalibrierungsbereichs ist nicht garantiert.
	
	\subsection{Residuale Oszillationen und die Nullstellen von \(\zeta(s)\)}
	
	Die Differenz zwischen den wahren Primzahlen und der theoretischen YUCT-Formel besteht aus deterministischen Oszillationen, die fast perfekt (\(R^2>0.99\)) durch die Beiträge der nichttrivialen Nullstellen der Riemannschen Zetafunktion beschrieben werden. Diese Beobachtung, die in Anhang PrimeN berichtet wird, legt nahe, dass diese Nullstellen die Eigenfrequenzen des Koordinationsfeldes sind. Eine vollständige parameterfreie geschlossene Form für \(p_n\) ohne jegliche Suche würde eine YUCT-basierte Herleitung der Nullstellen erfordern und bleibt zukünftiger Arbeit vorbehalten.
	
	\textbf{Status der Ergebnisse:}
	\begin{itemize}
		\item Theoretische YUCT-Formel: \textbf{Theorem} (asymptotisch).
		\item Empirische YUCT-Formel: \textbf{Empirisches Werkzeug}, kein Theorem.
		\item Spektrale Verbindung mit \(\zeta(s)\)-Nullstellen: \textbf{Empirische Beobachtung}, theoretische Erklärung offen.
	\end{itemize}
	
	%%%%%%%%%%%%%%%%%%%%%%%%%%%%%%%
	% TEIL IV: 19-DIMENSIONALE ELTERNTHEORIE
	%%%%%%%%%%%%%%%%%%%%%%%%%%%%%%%
	\part{19-dimensionale Elterntheorie}
	
	\section{Explizite Form der 19-dimensionalen Wirkung}
	\label{sec:19Dlagrangian}
	\[
	S_{19} = \frac{1}{2\kappa_{19}} \int d^{19}X \sqrt{-G}\; R(G) \;+\; \int d^{19}X \sqrt{-G}\; \mathcal{L}_{\text{coord}},
	\]
	mit
	\[
	\mathcal{L}_{\text{coord}} = -\frac{1}{4} \Tr(\Psi_{MN}\Psi^{MN}) + \frac{1}{2} G^{MN}\,\partial_M K_{\mathrm{eff}}\,\partial_N K_{\mathrm{eff}} - V(K_{\mathrm{eff}}) + \dots
	\]
	\textbf{Status:} Postulat.
	
	%%%%%%%%%%%%%%%%%%%%%%%%%%%%%%%
	% TEIL V: VORHERSAGEN UND PROGRAMM
	%%%%%%%%%%%%%%%%%%%%%%%%%%%%%%%
	\part{Testbare Vorhersagen und zukünftige Arbeit}
	
	\section{Konkrete experimentelle Tests}
	\begin{itemize}
		\item \textbf{Tensor-zu-Skalar-Verhältnis:} \(r \approx 0.07\).
		\item \textbf{Submikrometer-Gravitation:} \(\Delta g/g \sim 10^{-12}\) bei 10~nm.
		\item \textbf{CHSH-Ungleichung:} keine Verletzung über \(2\sqrt{2}\) hinaus.
		\item \textbf{Kosmisches Budget:} \(\Omega_\Lambda = 2/3\), \(\Omega_{\coord} = 0.283\), \(\Omega_b = 0.05\).
		\item \textbf{Temperaturabhängigkeit von \(G\):} \(\Delta G/G \sim 10^{-16}\).
		\item \textbf{Primzahl-asymptotische Skalierung:} der lokale Exponent \(\gamma\) muss für \(n\to\infty\) gegen \(2/3\) konvergieren (bereits für \(n\le 5\times10^5\) verifiziert).
	\end{itemize}
	
	\section{Programm für eine vollständige Herleitung}
	\begin{enumerate}
		\item Konstruktion der inneren Mannigfaltigkeit \(X_{15}\) und Berechnung der Fermionen-Massenparameter.
		\item Herleitung der effektiven Wirkung aus der 19-dimensionalen Theorie.
		\item Ab-initio-statistische Mechanik des Wörterbuchs.
		\item Herleitung der nichttrivialen Nullstellen von \(\zeta(s)\) aus den spektralen Eigenschaften des Koordinationsfeldes.
	\end{enumerate}
	
	\section{Fazit}
	Wir haben eine in sich geschlossene Herleitung des universellen Exponenten \(\beta = 2/3\) aus den Axiomen von YUCT präsentiert, einschließlich einer expliziten Kostenfunktionsanalyse und des Theorems des optimalen Transports. Es wird nun gezeigt, dass derselbe Exponent die asymptotische Verteilung der Primzahlen bestimmt, was die Theorie sowohl stärkt als auch eine neue Brücke zwischen Zahlentheorie und Koordinationsphysik öffnet. Die verbleibenden Elemente der Theorie sind klar als Vermutungen oder effektive Beschreibungen gekennzeichnet, was eine solide Grundlage für weitere Forschung bietet.
	
	\begin{center}
		\textit{Das Universum ist ein sich selbst koordinierendes Wörterbuch.}
	\end{center}
	
	\begin{thebibliography}{99}
		\bibitem{YakushevLaw2026}
		A.\,V.\,Yakushev, \emph{Yakushev's Law of Coordination: The Fundamental Law of Reality as a Fractal Coordination Network}, Zenodo (2026), DOI:10.5281/zenodo.18444598.
		\bibitem{AppendixL}
		A.\,V.\,Yakushev, ``Appendix L: Fractal Coordination Error Scaling — A Universal Law from DNA to Cosmology,'' in \emph{YUCT}, Zenodo (2026).
		\bibitem{AppendixY}
		A.\,V.\,Yakushev, ``Appendix Y: Mathematical Foundations of YUCT — A Derivation from First Principles,'' in \emph{YUCT}, Zenodo (2026).
		\bibitem{AppendixPrimeN}
		A.\,V.\,Yakushev, ``Appendix PrimeN: YUCT Prime Numbers Coordination Ladder,'' in \emph{YUCT}, Zenodo (2026).
		\bibitem{Rosser1941}
		J.\,B.\,Rosser, ``The \(n\)-th Prime,'' \emph{Amer. Math. Monthly} \textbf{48}, 607--611 (1941).
		\bibitem{West1997}
		West, G. B., Brown, J. H., \& Enquist, B. J. (1997). ``A general model for the origin of allometric scaling laws in biology.'' \emph{Science} \textbf{276}, 122--126.
		\bibitem{Banavar1999}
		Banavar, J. R., Maritan, A., \& Rinaldo, A. (1999). ``Size and form in efficient transportation networks.'' \emph{Nature} \textbf{399}, 130--132.
	\end{thebibliography}
	
\end{document}